\documentclass{article}
\usepackage{amsmath}
\begin{document}
\title{The useful observation has a recipient: relay placement and return in a cave projection}
\maketitle

\paragraph{Scope and input.}
The development family uses the published SubT Cave 01 labels
and authored unit-cost movement on its visibility graph.
The old report is sent from tile 18 to the fixed receiver 39
in slot zero under the unchanged native selected contact;
the old transmitter is not a third mobile egress obligation.
The report courier starts at 16 and the installer at 13;
both must egress at 0. The installer may use at most one
relay, at 13 or 42, through a named actual old-receipt cut.
Installation is effective in the slot after issue, and the
installer remains free to egress. Simultaneous independent
moves are permitted, with no collision or turning geometry.
A discovery move into a closed edge stays in place and uses
one slot. This is an admitted safe transition of the fixture,
not a claim that attempting a blocked passage is safe.
Selected closure timings and one omitted dispatch before the
broker are authored disturbances, not empirical probabilities.
The second dispatch remains owed until a positive receipt.

\paragraph{Native contact seam.}
For all 44 projected origins reachable by the courier in the
maximum horizon, run the pinned visibility, route registration,
deterministic 72 m hard range, broker dispatch and receiving
callback methods with relay absent, installed at 13, or at 42.
This gives 132 recorded native contact outcomes. The sampled
origin map and no-fade parameters are explicitly selected; actual
occlusion, propagation and forwarding back to the actor are
not supplied by these checks. The same contact table is used
throughout the model. Blocking movement edges does not rewrite
radio geometry, and no relative clock or reply path is inferred
from the broker. In particular, tile 6 is not in contact without
installation, while the contact from 42 can work without it.
A settled report is never spent twice.

\paragraph{Two exact policy comparisons.}
For a common-observation ideal, backward induction on time
enumerates both authorized actions at each attained state and
partitions the possible worlds by the resulting common receipt
and movement observations. It wins iff some action pair
wins on every resulting observation class. Present-closure
shortest-path pruning is valid because the modeled closures do
not reopen and a send occupies an additional slot. Reference
arms admit neither relay, only the near or far site, both
sites, only movement on already proved open edges, a separately
charged exact stationary scan, or complete knowledge of the world.
The stationary scan consumes one slot; it is a modeled alternative,
not a sensor proved to be available.

For two possible worlds, an actor's own chronological transcript
is either identical in both or has split the worlds permanently.
Until split, its two actions must coincide. After split, each
can be chosen independently with the common plan known.
The second solver carries both physical states and these two
knowledge bits. Thus it exactly enumerates local deterministic
strategies for the admitted communications and feedback.
It agrees with the common-observation solver when all observations
are shared. A routed one-bit signal can substitute when the
installer has already distinguished which world holds; the carrier
does not otherwise learn the installer's motion or the report
addressed only to that installer. Signal timing and authorization
are explicit ideal observation arms. Receiving the report at
the base does not itself produce those signals or the modeled
receipt feedback to the carrier.

\paragraph{Discriminating cases.}
In one pair either edge \(2\)--\(3\) or edge \(13\)--\(42\) is
closed and the latter world additionally suppresses slot-1
dispatch. The carrier's relay decision deadline is not its own
motion cut: at slot zero the installer attempts \(13\to42\)
while the carrier waits. At slot one, when the courier can learn
the installer outcome immediately and that installer has the
qualified earlier receipt, the courier chooses 42 or 6 while
the installer stages the far or near relay, respectively.
The source validates the slot-2 send in both branches and both
actors egress by 10. Without the installer-to-courier
distinction by slot one, there is no exact local solution at
this bound; sending it in the opposite direction or a slot
later does not suffice. An equally informed scheduler ties,
and a charged stationary scan at the installer also ties over
the sampled common-observation comparison. Where a separate
forced early attempt actually fails, the replay retains custody
and finishes the retry, rather than counting the planned send
as settlement.

Across 100 tested local-information configurations, private
observations suffice for 29; the ideal timely installer bit
suffices for 50, matching common observation. One additional
slot of latency reduces that count to 41; sharing only the
courier outcomes does not enlarge the 29 in this sweep.
These are counts in two-world directed information experiments,
not reliability estimates or general claims about antenna or
sensor performance. Earlier knowledge reduces a planning
obstruction but does not supply the signaling channel, make
physical traversal safe, or satisfy independent I2 acceptance.
\end{document}
